\small The search for long-lived particles provides an exciting and well motivated place that Beyond The Standard Model physics could be lurking. This dissertation presents contributions to two complementary ATLAS searches for charged long-lived particles with lifetimes of order 1~ns.

The first extends the ATLAS disappearing track search to previously unexplored \stau benchmark models, including the \stau coannihilation and GMSB scenarios. The observed results are consistent with Standard Model expectations within 2$\sigma$, leading to exclusion limits of 320~\GeV for a 1~ns lifetime in the CMSSM \stau coannihilation scenario and 300 \GeV for a 1~ns lifetime in the GMSB \stau scenario.

The second introduces a new search strategy combining disappearing track signatures with high ionization energy loss. The analysis employs a fully data-driven background estimation using two complementary methods: an ABCD approach and a template-based method. Validation studies demonstrate closure of both techniques, while projected sensitivity studies indicate significantly improved sensitivity to charged long-lived particles with lifetimes near 1~ns.

This dissertation also presents contributions to the development of the ATLAS Inner Tracker for HL-LHC operation, including studies of module performance under thermal cycling and the design of dedicated test infrastructure for loaded local support structures.

Together, these contributions advance both the detector capabilities and analysis techniques required for future charged long-lived particle searches at ATLAS.